\subsection{Certified native measurements}

The interval benchmark retains 420 measured calls and 84 warm-ups over
14 fixtures, with five shuffled rounds and separate A/A copies of the
historical, adaptive and pure-queue arms. All complete and return identical
certificates across arms. The 101.440-second execution includes fresh pairing
construction, complete counting, certificate generation and independent replay
in every timed call. Native meridian arc geometry is prepared outside these
interval timers. Serialization and cross-arm certificate comparisons are also
outside the timers. Table entries are separate median milliseconds; ratios
are medians of the five paired old/new ratios, not ratios of those medians.

\begin{center}\small
\begin{tabular}{@{}lrrrrr@{}}
\toprule
Input & old ms & adaptive ms & queue ms & old/adapt. & old/queue \\
\midrule
five-cycle-8 & 1.078 & 0.986 & 1.028 & 1.134 & 1.047 \\
five-cycle-16 & 5.567 & 3.046 & 3.324 & 1.828 & 1.687 \\
five-cycle-32 & 40.857 & 11.152 & 11.539 & 3.721 & 3.380 \\
five-cycle-64 & 303.463 & 42.773 & 45.284 & 7.119 & 6.772 \\
five-cycle-128 & 2405.376 & 167.526 & 179.342 & 14.413 & 13.389 \\
duplicates-16 & 0.292 & 0.308 & 1.113 & 0.969 & 0.257 \\
duplicates-64 & 1.051 & 1.056 & 16.320 & 0.991 & 0.064 \\
duplicates-256 & 4.090 & 4.107 & 291.486 & 1.003 & 0.014 \\
meridian-16 & 9.462 & 9.063 & 9.818 & 1.048 & 0.958 \\
meridian-64 & 241.530 & 192.101 & 203.347 & 1.290 & 1.188 \\
meridian-128 & 1578.288 & 1088.357 & 1163.742 & 1.463 & 1.364 \\
binary-64 & 7.847 & 3.714 & 3.876 & 2.037 & 2.043 \\
binary-1024 & 11.147 & 4.948 & 4.859 & 2.253 & 2.291 \\
binary-4096 & 24.273 & 9.387 & 8.568 & 2.587 & 2.828 \\
\bottomrule
\end{tabular}
\end{center}


At $m=128$, the five-cycle input takes 2,405.376 ms historically and
167.526 ms adaptively, with paired ratio 14.413. Merger tests fall from
2,105,660 to 97,409; the pure queue makes 81,154 tests but takes 179.342 ms.
The adaptive heap peaks at 11,907 entries, versus 12,097 for the pure queue.
The five cycles and 15,396-byte certificate are unchanged. This exhibits the
proved removal of redundant candidate tests while also showing that test
counts alone do not predict Python heap overhead.

For 256 identical rows, the old and adaptive engines each make 255 tests;
adaptive closure constructs no heap. Their times are 4.090 and 4.107 ms,
with paired ratio 1.003. Pure queue closure makes 65,025 tests, peaks at
48,769 heap entries and takes 291.486 ms. The adaptive first-pair fast path
therefore matters. At 16 duplicate rows the adaptive ratio is 0.969, a small
regression rather than a universal speedup claim.

For the supplied 128-tetrahedron meridian, the paired adaptive ratio is
1.463 (1,578.288 versus 1,088.357 ms). Both engines make 349,764 candidate
tests, and the adaptive engine never enters the queue: this gain comes from
indexing the periodic rows instead of rescanning nonperiodic rows inside each
inner loop. All 133 cycles and the 1,776,281-byte certificate agree. At fixed
$m=16$, increasing the binary scale to 4,096 bits gives adaptive ratio 2.587
and pure-queue ratio 2.828, with the same deterministic candidate counts as
the small-coordinate case.

Across interval fixtures the median A/A ratios range from 0.974 to 1.021
historically, 0.965 to 1.021 adaptively, and 0.956 to 1.013 for the queue.
Individual observations can vary more: the smallest five-cycle fixture
contains a 4.045-ms historical call alongside roughly 1-ms calls. All rounds
are retained; neither those observations nor warm-ups are silently discarded.
These are supplied interval systems, not whole-knot recognition timings.

The separate normal-query benchmark takes 539.262 seconds and retains
180 measured calls plus 36 warm-ups across nine fixtures, with historical
and adaptive arms and an A/A copy of each. Each timed call includes fresh
validation of the original triangulation and vector, geometry construction,
the requested census or reduced-core query, certificate production and the
independent normal-source verifier. Both engines answer exactly the same
query; every call completes and every certificate is identical across arms.
No tests, article builds or other benchmark drivers ran concurrently.

\begin{center}\small
\begin{tabular}{@{}lrrrrr@{}}
\toprule
Query & old ms & new ms & old/new & old A/A & new A/A \\
\midrule
disk-16 & 18.453 & 18.267 & 1.020 & 0.981 & 1.038 \\
disk-64 & 301.444 & 242.123 & 1.224 & 0.998 & 0.903 \\
disk-128 & 1691.973 & 1246.086 & 1.361 & 0.989 & 0.999 \\
coordinates-16 & 53.690 & 53.090 & 1.008 & 1.002 & 1.002 \\
coordinates-64 & 2224.639 & 2206.549 & 1.019 & 0.995 & 1.001 \\
coordinates-128 & 16871.366 & 16881.820 & 1.014 & 0.999 & 0.999 \\
core-64 & 21.728 & 21.252 & 1.030 & 0.999 & 1.012 \\
core-4096 & 23.322 & 22.730 & 1.021 & 1.007 & 0.993 \\
core-32768 & 35.037 & 34.245 & 1.025 & 0.993 & 0.999 \\
\bottomrule
\end{tabular}
\end{center}


For the 128-tetrahedron disk query, the complete paired ratio is 1.361:
1,691.973 ms historically versus 1,246.086 ms adaptively. The five individual
ratios lie between 1.265 and 1.441. There are still 133 orbit cycles,
21,008 weighted replay events and at most eight weight runs; the certificate
is still 1,777,349 bytes. The change is confined to finding the same periodic
merges more efficiently. All nine normal fixtures retain their merger-test
counts and take the indexed-scan path without switching to a heap.

The full-coordinate census retains much more work outside merger scheduling.
At 128 tetrahedra it carries 896 coordinates and reaches 385 weight runs;
its unchanged certificate has 2,504,362 bytes. Separate old/new medians are
16,871.366 and 16,881.820 ms. Although the paired median ratio is 1.014,
the five ratios are 1.014, 1.021, 1.026, 0.629 and 0.733. In the final two
rounds the new arm takes 26,811.822 and 23,415.500 ms, well above its earlier
roughly 16--17-second calls. The median A/A ratios near one do not explain
away this variation. These data establish no reliable full-coordinate
speedup. The 64-tetrahedron disk query also has a new-arm A/A ratio of 0.903;
its paired ratios range from 0.914 to 1.434, so its median 1.224 gain has
appreciable noise.

The core fixtures use $x=gD+(g+1)L$, $g=2^b$, on the 16-tetrahedron solid
torus, with $D$ its meridian vector and $L$ its vertex-link vector. Both
engines already apply the same quadrilateral normalization. Their paired
ratios range only from 1.021 to 1.030; at $b=32,768$ the medians are 35.037
and 34.245 ms. The unchanged reduced query still has 21 cycles, 400 events
and a 56,757-byte source-bound certificate. These small timing differences
are not a second claim of the much larger normalization gain reported in
the preceding weighted-component section. Across normal fixtures the
historical median A/A ratios range from 0.981 to 1.007 and the new ratios
from 0.903 to 1.038. All 600 measured calls and 120 warm-ups in the two
new benchmarks complete, but none is a whole-knot recognition measurement.
